\documentclass[12pt,a4paper]{article}
\usepackage[utf8]{vietnam}
\usepackage{geometry}
\geometry{left=2.5cm,right=2.5cm,top=2.5cm,bottom=2.5cm}
\usepackage{graphicx}
\usepackage{listings}
\usepackage{xcolor}
\usepackage{hyperref}
\hypersetup{
    colorlinks=true,
    linkcolor=blue,
    urlcolor=blue,
}
\definecolor{codebg}{rgb}{0.95,0.95,0.95}
\lstset{
    backgroundcolor=\color{codebg},
    basicstyle=\ttfamily\small,
    breaklines=true,
    frame=single,
    language=bash,
    showstringspaces=false,
}

\title{\textbf{Get Your Head In The Cloud with Google Cloud Platform} \\ \textit{Practical Session - Building a Serverless Image Search Engine on GCP}}
% \author{Nguyen Hoang Long \\ \texttt{longnguyenhoang.personal@gmail.com}}
\date{}

\begin{document}

\maketitle

\tableofcontents
\newpage

\section{Introduction}
In this hands‑on session, you will build a simple but complete \textbf{image search engine}. Users can upload images with descriptive keywords, and later search for images using those keywords. The application runs entirely on Google Cloud Platform (GCP) using two core services:

\begin{itemize}
    \item \textbf{Cloud Storage} – stores uploaded images and a metadata file that maps keywords to image URLs.
    \item \textbf{Cloud Run} – runs the web application (written in Python/Flask) without you having to manage any servers.
\end{itemize}

The solution is \textbf{serverless}: it scales automatically from zero to many instances based on traffic, and you pay only for what you use. By the end of this session, you will have deployed a real, publicly accessible application on GCP.

\section{Prerequisites}
Before starting, ensure you have:
\begin{itemize}
    \item A Google account.
    \item A GCP project with billing enabled (you can use the \$300 free trial).
    \item Access to the \href{https://console.cloud.google.com}{Google Cloud Console}.
\end{itemize}
All commands will be run from \textbf{Cloud Shell}, which is a free Linux environment built into the console. It already has the \texttt{gcloud} CLI and other tools pre‑installed.

\section{Step-by-step implementation}

\subsection*{Step 1: Open \texttt{Cloud Shell} and enable APIs}
\begin{enumerate}
    \item Go to \url{https://console.cloud.google.com}.
    \item Make sure you are in the correct project (create one if needed).
    \item Click the \texttt{>\_} icon in the top right corner to open Cloud Shell. A terminal will appear at the bottom of the page.
    \item Enable the required APIs:
\begin{lstlisting}
gcloud services enable storage.googleapis.com run.googleapis.com
\end{lstlisting}
    This command allows your project to use Cloud Storage and Cloud Run.
\end{enumerate}

\subsection*{Step 2: Create a Cloud Storage Bucket}
A bucket is a container for your files (images and metadata). The name must be globally unique, so we will append a timestamp.
\begin{lstlisting}
BUCKET_NAME="image-search-$(date +%s)"
gsutil mb gs://$BUCKET_NAME
\end{lstlisting}
Verify the bucket was created:
\begin{lstlisting}
gsutil ls
\end{lstlisting}
You should see your bucket name listed.

\subsection*{Step 3: Prepare the application code}
We will create a new directory and write a Python web application using Flask.

\begin{lstlisting}
mkdir image-search && cd image-search
\end{lstlisting}

Create the main application file \texttt{app.py}:

\textbf{Explanation of the code:}
\begin{itemize}
    \item The app uses Flask to serve two main endpoints:
        \begin{itemize}
            \item \texttt{/} – displays the upload form and search results.
            \item \texttt{/upload} – handles image uploads, stores the image in Cloud Storage, and updates the metadata file.
            \item \texttt{/search} – reads the metadata and returns matching images.
        \end{itemize}
    \item Metadata is stored in a JSON file (\texttt{keywords.json}) inside the bucket. Each keyword maps to a list of image URLs and associated information.
    \item Cloud Storage is accessed via the \texttt{google-cloud-storage} client library.
\end{itemize}

Now create the \texttt{requirements.txt} file:
\begin{lstlisting}
Flask==2.3.3
google-cloud-storage==2.9.0
gunicorn==21.2.0
\end{lstlisting}

Finally, create a \texttt{Dockerfile} to define the container environment:
\begin{lstlisting}
FROM python:3.9-slim
WORKDIR /app
COPY requirements.txt .
RUN pip install --no-cache-dir -r requirements.txt
COPY . .
CMD exec gunicorn --bind :$PORT --workers 1 --threads 8 app:app
\end{lstlisting}

\subsection*{Step 4: Deploy to Cloud Run}
Now we will deploy the application. Cloud Run will automatically build the container from your source code and run it.
\begin{lstlisting}
gcloud run deploy image-search \
  --source . \
  --platform managed \
  --region asia-southeast1 \
  --allow-unauthenticated \
  --set-env-vars BUCKET_NAME=$BUCKET_NAME
\end{lstlisting}

\textbf{What each flag does:}
\begin{itemize}
    \item \texttt{--source .} – tells Cloud Run to build the image from the current directory.
    \item \texttt{--platform managed} – uses the fully managed Cloud Run environment (no clusters to manage).
    \item \texttt{--region asia-southeast1} – deploys the service in Singapore, which is close to Vietnam for low latency.
    \item \texttt{--allow-unauthenticated} – makes the service publicly accessible (no login required).
    \item \texttt{--set-env-vars} – passes the bucket name to the application as an environment variable.
\end{itemize}

After a few moments, the command will output a URL, e.g., \texttt{https://image-search-xxxxx-uc.a.run.app}. Open this URL in your browser.

\subsection*{Step 5: Test your application}
\begin{enumerate}
    \item You should see a page with an upload form and a search box.
    \item Upload an image. Enter some keywords separated by commas, e.g., \texttt{meme, funny}.
    \item After uploading, you will be redirected back to the home page.
    \item Try searching for one of the keywords you used. You should see the image(s) you uploaded.
    \item Upload another image with different keywords and verify that searching works.
\end{enumerate}

You can also explore the bucket in the Cloud Console:
\begin{itemize}
    \item Go to \textbf{Cloud Storage} → \textbf{Buckets}.
    \item Click on your bucket name.
    \item You will see a folder \texttt{images/} containing the uploaded files, and a file \texttt{keywords.json} that holds the search index.
\end{itemize}

\subsection*{Step 6: Clean up (important)}
To avoid incurring charges after the session, delete the resources you created.
\begin{lstlisting}
# Delete all files in the bucket and then the bucket itself
gsutil rm -r gs://$BUCKET_NAME

# Delete the Cloud Run service
gcloud run services delete image-search --region asia-southeast1 --quiet
\end{lstlisting}
If you no longer need the project, you can delete it from the Cloud Console.

\section{What You Learned}
\begin{itemize}
    \item How to create a Cloud Storage bucket and upload files.
    \item How to deploy a serverless web application using Cloud Run.
    \item How to use Python and Flask to interact with Cloud Storage.
    \item A practical example of separating data (images, metadata) from application logic.
\end{itemize}

\section{Possible Enhancements}
If you want to explore further, consider these extensions:
\begin{itemize}
    \item Add authentication so only authorized users can upload.
    \item Use Firestore instead of a JSON file for better concurrency and scalability.
    \item Automatically generate thumbnail images using Cloud Functions.
    \item Add a simple voting system for memes.
\end{itemize}

\vspace{1cm}
\begin{center}
\textbf{Arigathanks for participating!} \\
Happy cloud building!
\end{center}

\end{document}